\newpage
\section{Introduction}
% ──────────────────────────────────────────────────────────────────────────────

Explainable Artificial Intelligence (XAI) refers to a set of methods that aim
to make machine learning models understandable to humans by revealing how their
internal mechanisms work and how decisions are reached~\cite{Arrieta2020XAI}. In medicine, this
capability is not merely desirable but increasingly required: clinical
decision-making carries direct consequences for patient safety, and regulatory
frameworks in many jurisdictions now demand that AI-assisted diagnoses be
accompanied by interpretable justifications~\cite{jin2023guidelines}. Despite
this pressure, metrics for measuring the effectiveness of XAI algorithms in
medical contexts remain poorly established. Current evaluation methodologies
rely heavily on subjective user studies which, while valuable for understanding
clinician perception, often suffer from confirmation bias and fail to assess the
technical fidelity of the explanations provided~\cite{mirzaei2023xai}.

Moreover, many modern XAI approaches are post-hoc in nature, which makes
evaluating their explanations even more challenging. Because such explanations
are generated after model training, they may not faithfully reflect the model's
true decision-making logic. This problem is especially acute in medicine, where
a persuasive but unfaithful explanation could reinforce a clinician's incorrect
judgement rather than correct it.

This paper focuses specifically on the medical domain and asks: which
quantitative evaluation metrics are actually being used, how do they
co-occur across studies, and where does coverage remain thin? By analysing a
focused corpus of peer-reviewed medical XAI studies across four sub-domains, including
oncology, cardiology, neurology, and radiology, we identify redundancies in
current practice and propose a minimal but comprehensive evaluation set grounded
in a three-layer taxonomy of mathematical faithfulness, system robustness, and
human-centred trust.


\section{Background}
% ──────────────────────────────────────────────────────────────────────────────

Artificial intelligence is rapidly being integrated into clinical workflows.
Yet because complex AI models often operate as a ``black box,'' it is difficult
for clinicians to understand how diagnostic or prognostic decisions are actually
being made. This creates a significant need for XAI, which aims to make medical
AI systems transparent and interpretable~\cite{Tjoa2021MedicalXAI}. XAI has been applied to problems such
as clinical decision support~\cite{abbas2025explainable}, breast cancer
diagnosis~\cite{saharan2025explainable}, cardiovascular risk
prediction~\cite{shah2025explainable}, and responsible cardiovascular disease diagnosis~\cite{muneer2025blockchain}.

A central distinction in XAI is between models that are interpretable by
design and methods that explain a model after the fact. The latter,
so-called post-hoc approaches, are appealing because they can be applied to
virtually any existing clinical system without requiring it to be rebuilt.
However, this flexibility comes at a cost: there is no guarantee that the
explanation accurately reflects what the model is doing internally, and this
gap between appearance and reality is one of the core challenges medical XAI
continues to grapple with~\cite{holzinger2022information,palatnik2021xai}.

Evaluating explanation quality in medicine has accordingly become a major
concern. Two broad approaches have emerged. \emph{Technical evaluation} uses
quantitative metrics, such as faithfulness, sensitivity, robustness, and
consistency, to assess how well an explanation approximates actual model
behaviour. \emph{Human-centric evaluation}, by contrast, focuses on the
clinician, measuring factors such as trust, cognitive load, and practical
usefulness through observational studies and behavioural experiments~\cite{lenis2020domain}.

While both paradigms have matured considerably, they remain largely
disconnected in medical practice. A method that scores well on technical
metrics may produce explanations that are opaque or misleading to clinicians,
while an explanation that non-specialists find intuitive may not be
mathematically faithful to the underlying model. This disconnect makes it
difficult to establish shared evaluation standards or compare results across
medical sub-domains, a problem this paper directly addresses.

% ──────────────────────────────────────────────────────────────────────────────